\section{An elementary Heisenberg carrier estimate}
\label{sec:heisenberg-carrier}

\paragraph{Status.}
This is a complete new elementary derivation. The author, parent,
and an independent sibling reviewer have read the entire proof and
passed it in the stated scope. Their exact checks are recorded in
CHECK\_EVENTS.tsv. No Weil-character formula or representation-degree
classification is a premise.
The statements concern ordinary complex representations.

\begin{lemma}[D-HOS-01: fixed-point-free symplectic carriers]
\label{D-HOS-01}
Let $q=p^f$ be odd, let $(W,B)$ be a symplectic
$\mathbf F_q$-space of dimension $2m\ge2$, and let
\[
 H=W\times\mathbf F_q,\qquad
 (v,z)(w,t)=(v+w,z+t+\tfrac12 B(v,w))
                                                               \tag{HOS.1}
\]
be its Heisenberg group. Let $g\in\operatorname{Sp}(W)$ have
prime order $\ell\ne p$ and satisfy $W^g=0$. Suppose that a
finite-order operator $T$ on an $a$-dimensional complex
$H$-module $V$ satisfies
\[
 T\rho(v,z)T^{-1}=\rho(gv,z).
                                                               \tag{HOS.2}
\]
Put $t=\dim V^H$ and
$\nu(T)=a-\max_\lambda\dim\ker(T-\lambda I)$. Then
\[
 \nu(T)\ge
 \left(1-\frac1\ell\right)(1-q^{-m})(a-t).
                                                               \tag{HOS.3}
\]
If moreover $|\operatorname{Tr}\rho(h)|\le ca$ for every
$h\in H\setminus\{1\}$, with $0\le c\le1$, then
\[
 \nu(T)\ge
 \left(1-\frac1\ell\right)(1-q^{-m})
       (1-q^{-(2m+1)})(1-c)a.
                                                               \tag{HOS.4}
\]
Neither $\rho$ nor its restriction to $H$ is assumed irreducible,
and $T^\ell$ is not required to be scalar on $V$.
\end{lemma}

\begin{proof}
We give the carrier construction and every multiplicity step.
Choose a symplectic decomposition $W=X\oplus Y$ identifying
$X$ and $Y$ with $\mathbf F_q^m$, so
\[
 B((x,y),(x',y'))=x\cdot y'-x'\cdot y.
\]
Fix a nontrivial additive character $\psi:\mathbf F_q\to\mathbf C^*$.
On the $N=q^m$-dimensional function space
$U_\psi=\{f:\mathbf F_q^m\to\mathbf C\}$ define
\[
 (L_{x,y}f)(u)
   =\psi(y\cdot u+\tfrac12 x\cdot y)f(u+x),\qquad
 R_\psi((x,y),z)=\psi(z)L_{x,y}.
                                                               \tag{HOS.5}
\]
Direct multiplication gives
$L_vL_w=\psi(B(v,w)/2)L_{v+w}$, so (HOS.5) is a
representation of (HOS.1). Its trace is zero at $L_v$ for
$v\ne0$: a nonzero translation has no diagonal entries,
whereas for $x=0,y\ne0$ the sum of the nontrivial additive
character $u\mapsto\psi(y\cdot u)$ is zero.
It follows that
\[
        \operatorname{Tr}(L_vL_w^{-1})=N\delta_{v,w}.
                                                               \tag{HOS.6}
\]
Thus the $q^{2m}=N^2$ operators $L_v$ form a basis of
$\operatorname{End}_{\mathbf C}(U_\psi)$, and this
representation is irreducible. It is the only irreducible
representation with central character $\psi$: the corresponding
central-character summand of $\mathbf C[H]$ has dimension
$|H|/q=q^{2m}$ and maps onto this full matrix algebra of
the same dimension. The other irreducibles, with trivial
central character, are exactly the linear characters of
$H/Z(H)=W$.

The automorphism $(v,z)\mapsto(gv,z)$ preserves $\psi$.
Uniqueness therefore supplies an intertwiner $\Omega$ on
$U_\psi$ such that
\[
                  \Omega L_v\Omega^{-1}=L_{gv}.
                                                               \tag{HOS.7}
\]
Schur's lemma makes $\Omega^\ell$ scalar; rescale $\Omega$
so that $\Omega^\ell=I$. Equation (HOS.7) is exact, not
an equality with additional phases, because it implements
the specified automorphism on the actual group elements
$(v,0)$ in (HOS.1).
For $1\le j<\ell$, the permutation of the basis
$\{L_v\}$ induced by conjugation with $\Omega^j$ has
exactly one fixed vector, $L_0$. Indeed, $W^{g^j}=W^g=0$
because $\ell$ is prime. Hence
\[
 |\operatorname{Tr}\Omega^j|^2
  =\operatorname{Tr}(\operatorname{Ad}\Omega^j)
  =|W^{g^j}|=1.
                                                               \tag{HOS.8}
\]
The first equality uses finite order: the eigenvalues of
$\Omega$ are roots of unity, so the trace of its inverse
is the complex conjugate of its trace. Fourier inversion
on the cyclic group of order $\ell$ now bounds every
eigenvalue multiplicity of $\Omega$ by
\[
                  \frac{N+\ell-1}{\ell}.
                                                               \tag{HOS.9}
\]
Its projective support is consequently at least
$(1-1/\ell)(1-N^{-1})N$.

Decompose $V$ into its $H$-isotypic spaces. Every summand
with nontrivial central character has the form
$U_\psi\otimes M_\psi$. Since $g$ fixes the center,
it is $T$-stable, and (HOS.7) plus the full matrix-algebra
commutant gives
\[
                 T=\Omega\otimes A_\psi
                   \quad\hbox{on }U_\psi\otimes M_\psi.
                                                               \tag{HOS.10}
\]
The operator $A_\psi$ has finite order after the indicated
normalization: some power divisible by $\ell$ kills both
$T$ and $\Omega$. It is therefore diagonalizable. The
normalized eigenvalue multiplicities of a tensor product
are convolutions, so their maximum is no greater than
the maximum for $\Omega$. The support fraction from
(HOS.9) holds on every such isotypic summand, with no
assumption about the eigenvalues on $M_\psi$.

For the linear characters, $g$ has no fixed nonzero dual
weight because $g-1$ is invertible on $W$. Each nonzero
dual orbit therefore has length $\ell$, with all weight
spaces in it of equal multiplicity. On its direct sum,
$T$ is a cyclic block operator. Its $\ell$th power has
the same spectrum on all blocks, and each eigenvalue
of the block product gives at most one copy of any
chosen eigenvalue of $T$ per multiplicity. Equivalently,
conjugate the cyclic block operator to one with identity
intertwiners except for a final diagonal block.
Every eigenspace has dimension at most $1/\ell$ of
that orbit summand. This support fraction $1-1/\ell$
is at least the one in (HOS.9).

The remaining trivial $H$-isotypic space has dimension
$t$ and is $T$-stable. Give it the worst possible full
eigenspace of dimension $t$. Adding the preceding
eigenspace upper bounds over all other $T$-stable
summands proves (HOS.3). Finally, the averaging projector
onto $V^H$ gives
\[
 t=\frac1{|H|}\sum_{h\in H}\operatorname{Tr}\rho(h)
       \le a\left(c+\frac{1-c}{q^{2m+1}}\right).
                                                               \tag{HOS.11}
\]
Substitution into (HOS.3) proves (HOS.4).
\end{proof}

\begin{lemma}[D-HOS-02: an elementary symplectic degree bound]
\label{D-HOS-02}
Let $q$ be odd and $n\ge4$. Every nontrivial ordinary irreducible
representation of $\operatorname{Sp}_{2n}(q)$ has degree
\[
                     a\ge\frac{q-1}{2}\,q^{n-1}.
                                                               \tag{HOS.12}
\]
\end{lemma}

\begin{proof}
The line-parabolic unipotent radical is the Heisenberg group
in (HOS.1) with $m=n-1$. Its center is the additive group
$\mathbf F_q$ of long-root unipotents. The kernel of a
nontrivial irreducible representation of the quasisimple
symplectic group is central, so this noncentral unipotent
group cannot act trivially. At least one nonzero additive
central character therefore occurs. By the explicit group
algebra argument in (HOS.5)--(HOS.6), its isotypic space
has dimension at least $q^{n-1}$.

The line-parabolic torus sending a chosen isotropic vector
to $s$ times itself and its paired dual vector to $s^{-1}$
times itself acts on the center by multiplication by $s^2$.
Parametrize the additive characters by
$\psi_w(z)=\exp(2\pi i\operatorname{Tr}_{q/p}(wz)/p)$.
The orbit of each $w\ne0$ under all squares in
$\mathbf F_q^\times$ has size $(q-1)/2$. Every character
in this orbit must occur, with the same isotypic dimension.
Summing these dimensions proves (HOS.12). The symplectic
quasisimplicity used here is the standard ordinary-group
input; no exceptional cover occurs in the stated range
by the complete list in \cite[Chapter~5, Table~1, p.~35]{BI:HHM}.
\end{proof}
